\section{Repartition des taches}

% Tableau indiquant qui a travaille sur quoi et dans quelle proportion.
% La repartition doit etre EQUILIBREE SUR LE FOND : aucun membre ne peut se
% cantonner a la redaction du rapport ou a la preparation de la presentation.
%
% A tenir a jour en continu, en s'appuyant sur l'historique git. Point de depart :
% `python report/contributions.py --latex` (lignes par domaine ; a valider a trois).

\begin{table}[h]
\centering
\begin{tabular}{lccc}
\toprule
\textbf{Tache} & \textbf{Hamza} & \textbf{Shaleena} & \textbf{Lenny} \\
\midrule
Modelisation mathematique      & & & \\
Schema numerique               & & & \\
Moteur et projection           & & & \\
Champ de guidage               & & & \\
Scenario et geometrie          & & & \\
Observables et campagnes       & & & \\
Validation                     & & & \\
Redaction du rapport           & & & \\
Transparents                   & & & \\
\bottomrule
\end{tabular}
\caption{Repartition des taches au sein du groupe, en pourcentage.}
\end{table}
